% 图契约：反应性近场的往返能流与辐射远场的净外流；几何、方向和公式与正文同一约定。
\documentclass[tikz,border=5pt]{standalone}
\usepackage{fontspec,xeCJK,amsmath,amssymb}
\IfFontExistsTF{Times New Roman}{\setmainfont{Times New Roman}}{\setmainfont{TeX Gyre Termes}}
\IfFontExistsTF{Heiti SC}{\setCJKmainfont[AutoFakeBold=1.5]{Heiti SC}}{\setCJKmainfont[AutoFakeBold=1.5]{FandolHei-Regular}}
\usetikzlibrary{arrows.meta,calc,patterns}
\definecolor{ink}{HTML}{243745}
\definecolor{blue}{HTML}{275DAB}
\definecolor{green}{HTML}{177C63}
\definecolor{red}{HTML}{BD4935}
\definecolor{orange}{HTML}{AA6C16}
\definecolor{purple}{HTML}{715393}
\begin{document}
\begin{tikzpicture}[x=1cm,y=1cm,font=\fontsize{9.5}{12}\selectfont,
v/.style={-{Latex[length=2mm,width=1.3mm]},line width=.85pt},
axis/.style={-{Latex[length=1.7mm,width=1.1mm]},ink!65,line width=.55pt},
aux/.style={ink!45,dashed,line width=.55pt},
lab/.style={fill=white,inner sep=1.8pt},
card/.style={draw=ink!25,fill=blue!3,rounded corners=2pt,text width=6.1cm,minimum height=.85cm,inner sep=6pt,align=center}]
\path[use as bounding box](-8,-3.6)rectangle(8,3.65);
\node[font=\bfseries\fontsize{10.5}{13}\selectfont,text width=15.1cm,align=center]at(0.00000,3.22000){M7-H51　反应性近场的往返能流与辐射远场的净外流};
\begin{scope}[shift={(-4,0)}]
\draw[ink,fill=blue!5](-1.5,0)circle(0.55);
\node[]at(-1.50000,1.50000){近场};
\draw[ink,fill=blue!5](1.5,0)circle(0.55);
\node[]at(1.50000,1.50000){远场};
\draw[{Latex[length=1.6mm]}-{Latex[length=1.6mm]},blue](-2.60000,0.00000)--(-2.10000,0.00000);
\draw[{Latex[length=1.6mm]}-{Latex[length=1.6mm]},blue](-0.90000,0.00000)--(-0.40000,0.00000);
\node[]at(-1.50000,-1.30000){能量随时间往返};
\draw[v,red](0.90000,0.00000)--(0.40000,0.00000);
\draw[v,red](2.10000,0.00000)--(2.70000,0.00000);
\node[]at(1.50000,-1.30000){平均能流净向外};
\end{scope}
\node[card]at(4.00000,1.25000){$\text{近场：储能往返交换}$};
\node[card]at(4.00000,-1.05000){$\text{远场：平均能流净向外}$};
\end{tikzpicture}\end{document}
